\documentclass[11pt]{article}

\usepackage[margin=1in]{geometry}
\usepackage{amsmath,amssymb}
\usepackage{booktabs}
\usepackage[hidelinks]{hyperref}
\usepackage{url}
\usepackage{microtype}

\title{Methods for the diagnostic principal-component analysis\\ of a six-asset technical-indicator matrix}
\author{}
\date{17 November 2026}

\newcommand{\tr}{\mathsf{tr}}

\begin{document}
\maketitle

\begin{abstract}
This note records the statistical methods that have been fitted in the \texttt{AI-trades/} experiment and no others.
Two representations of the same \(6{,}025\times 228\) indicator matrix are retained: the stored columns, and a price-relative transformation of the \(132\) columns whose units are prices or price differences.
Each representation is standardized and decomposed on rows through 31~December~2018 only.
Component-retention rules, including a classical Marchenko--Pastur edge, are compared and are not selected.
Section~\ref{sec:roadmap} records the proposed Capstone LDI roadmap.
Stages~3--7 of that roadmap are not fitted methods.
\end{abstract}

\section{Scope}

The research program writes \path{data/asset_indicator_matrix.csv} from \texttt{main.py}.
The experiment reads that file and does not import the research package.
A method appears below only when the corresponding fit or check has been run and saved.
Quantities that were computed and then rejected as decision rules are identified as such.
The checkpoint is the second diagnostic decomposition, stored as \texttt{stage2\_price\_relative\_train2018\_v1}.
Sections from Section~\ref{sec:data} onward describe that fitted work.
They do not implement the later roadmap stages.

\section{Research Architecture, Milestones, and Deliverables}
\label{sec:roadmap}

This section is the Capstone LDI research roadmap.
Roadmap Stages~0 and~1 are complete.
Later stages are plans, and updating this note does not accept them.
The saved model directories use an earlier internal numbering.
\path{models/stage1_train2018_v1/} is the raw-matrix diagnostic PCA.
\path{models/stage2_price_relative_train2018_v1/} is the price-relative diagnostic PCA.
Under the numbering below, those two fits are evidence inside Stages~1 and~2.
Roadmap Stage~1 is closed by the approval recorded in that stage.
Roadmap Stage~2 is not closed.
Writing this section does not approve Stage~3.

\subsection{Stage-completion policy}

A stage is complete only after its artifacts exist, validation passes, documentation is updated, and the project owner explicitly approves it.
Stage~0 retains the complete status recorded when this roadmap was added.
From Stage~1 onward, that status is not changed to complete by writing artifacts or by updating this note.

\subsection{Stage 0 --- Feature Engineering}

\textbf{Objective.} Construct the historical multi-asset technical-indicator matrix.

\textbf{Deliverable.} The unstandardized file \path{asset_indicator_matrix.csv}.

\textbf{Status.} Complete.
That file is the research matrix described in Section~\ref{sec:data}.
This experiment reads it and does not write it.

\subsection{Stage 1 --- Data Preparation and Statistical Validation}

\textbf{Objective.} Audit feature definitions, normalize price-denominated indicators, and establish leakage-safe standardization.

\textbf{Deliverables.} A reusable transformed and standardized feature matrix, the fitted preprocessing pipeline, transformation metadata, and validation diagnostics.

\textbf{Status.} Complete / Approved on 17~November~2026.
The approved artifact is \path{AI-trades/artifacts/stage1_price_relative_standardized_v2/}.
The four completion gates are satisfied: the artifacts exist, the validation checks recorded below passed, this note records them, and the project owner has approved the stage.
No column count or retention rule was selected as part of Stage~1.

The export is a package of the price-relative map in Section~\ref{sec:transform} and the frozen scaler
\begin{center}
\path{models/stage2_price_relative_train2018_v1/scaler.joblib}.
\end{center}
That scaler was not refit.
The directory is \path{AI-trades/artifacts/stage1_price_relative_standardized_v1/}, created 2026-11-17T21:55:47Z.
It contains \path{feature_matrix_standardized.csv}, \path{preprocessing_pipeline.joblib}, \path{feature_metadata.json}, and \path{validation_report.md}.
The matrix has \(6{,}025\) rows and \(228\) columns.
Its columns keep the source MultiIndex: the CSV has an asset header row, an indicator header row, and a date index, and it is read back with \texttt{header=[0,1]} and \texttt{index\_col=0}.
The \(4{,}044\) training rows have maximum absolute column mean \(2.09\times 10^{-14}\) and maximum absolute gap from unit population variance \(4.44\times 10^{-16}\).
The other \(1{,}981\) rows were transformed with the same frozen scaler.
That scaler records \(4{,}044\) samples, and the largest absolute test-column mean is \(0.746\), so the full sample was not used to standardize.
The training block matches the Stage~2 standardized PCA input with maximum absolute error \(0\) against the saved scaler's affine map and against the saved PCA training mean.
Reloading the pipeline and applying it to the same inputs reproduces the exported matrix with maximum absolute error \(0\).
The CSV round trip differs by at most \(3.55\times 10^{-15}\).
The source hash and adjusted-close hash match the Stage~2 metadata.
The numerical record is \path{artifacts/stage1_price_relative_standardized_v1/validation_report.md}.
That v1 joblib file is historical evidence.
It names the class \texttt{\_\_main\_\_.FrozenPriceRelativePreprocessor}.
That name does not load in a fresh interpreter.
The portable Stage~1 artifact is
\path{artifacts/stage1_price_relative_standardized_v2/},
created 2026-11-17T22:10:39Z.

Its class is \texttt{ai\_trades.preprocessing.FrozenPriceRelativePreprocessor}.
The v2 standardized matrix is a byte-for-byte copy of the v1 matrix.
The scaler was not refit and still records \(4{,}044\) observations.
A fresh process loaded v2 with ordinary \texttt{joblib.load}.
On 2026-11-04 through 2026-11-17 the resulting \(10\times 228\) matrix differed from v1 by at most \(8.882\times 10^{-16}\), and a repeated call differed by zero.
The v2 record is \path{artifacts/stage1_price_relative_standardized_v2/validation_report.md}.

\subsection{Stage 2 --- PCA and Random Matrix Analysis}

\textbf{Objective.} Extract principal components and study the eigenvalue spectrum, redundancy, and candidate retention rules.

\textbf{Deliverables.} A time-indexed PCA score matrix, loadings, eigenvalues, explained variance, retention candidates, and diagnostics.

\textbf{Status.} Pending --- awaiting the next research session.
That session is to begin with a conceptual and mathematical review of principal components before any new implementation.
The diagnostic fits already saved are described below.
They are not the completed Stage~2 deliverable, and approval of Stage~1 does not start Stage~2.
Both saved decompositions report eigenvalues, explained-variance comparisons, retention candidates, and diagnostics.
The reports are \path{reports/stage1_diagnostics_v1.md} and \path{reports/stage2_diagnostics_v1.md}.
A reusable time-indexed score matrix has not been exported as a stage deliverable, and no retention rule has been selected.
The Marchenko--Pastur calculations in Section~\ref{sec:mp} remain diagnostics.

\subsection{Stage 3 --- Market Regime Identification}

\textbf{Objective.} Identify statistically stable, recurring market regimes using unsupervised clustering.

\textbf{Deliverables.} Historical regime classifications, a saved clustering model, stability diagnostics, and regime interpretation.

\textbf{Status.} Planned.
No clustering model has been fit.

\subsection{Stage 4 --- Markov Regime Dynamics}

\textbf{Objective.} Estimate regime transitions and multi-session regime probabilities.

\textbf{Deliverables.} A transition matrix, five- and ten-session regime distributions, and transition-model diagnostics.

\textbf{Status.} Planned.
No transition matrix has been estimated.

\subsection{Stage 5 --- Conditional Return Modeling}

\textbf{Objective.} Estimate and evaluate regime-conditioned QQQ return distributions.

\textbf{Deliverables.} Five- and ten-session return forecasts, uncertainty estimates, and out-of-sample validation.

\textbf{Status.} Planned.
No conditional return distribution has been estimated.
Test-period profitability is not available as a criterion for choices in Stages~1 or~2.

\subsection{Stage 6 --- Strategy Validation}

\textbf{Objective.} Translate forecasts into testable trading rules and evaluate performance.

\textbf{Deliverables.} A reproducible backtest, transaction-cost assumptions, benchmark comparisons, drawdown analysis, and a performance report.

\textbf{Status.} Planned.
No trading rule or backtest has been run.

\subsection{Stage 7 --- Paper Trading}

\textbf{Objective.} Evaluate the strategy under realistic forward-looking conditions.

\textbf{Deliverables.} Daily signal reports, a dated paper-trading log of those signals, and a standing constraint that a report does not submit an order.
Until a validated rule exists, the default signal is no trade.
A daily report would apply a previously fitted model to the latest row and would not refit that model.

\textbf{Status.} Planned.
No signal report or paper-trading log has been produced.

\section{Data and timing conventions}
\label{sec:data}

Let \(t\) index exchange sessions and let
\[
\mathcal{A}=\{\mathrm{SPY},\mathrm{QQQ},\mathrm{IWM},\mathrm{TLT},\texttt{\^{}VIX},\texttt{DX-Y.NYB}\}
\]
be the asset universe.
The observed matrix \(X\) has one row per session and columns \((a,j)\) for \(a\in\mathcal{A}\) and \(j\) in a set of \(38\) indicators, so \(p=6\times 38=228\).
The archived copy used by both fits is \path{AI-trades/data/feature_matrix_v1.csv}:
\(n=6{,}025\) rows, from 4~December~2002 through 17~November~2026, SHA-256
\begin{center}
\path{4ecefa23bb0321deea57d55a05bc45cba11424ba852554ae40b1e3666ed457ad}.
\end{center}
The matrix contains no missing values.
A later hash of the research file is copied to the next \texttt{feature\_matrix\_v\(N\).csv} and appended to \path{data/manifest.jsonl}.
An existing snapshot is not replaced.
No second feature-matrix version has been produced.

Timestamps are reduced to a timezone-naive calendar date by discarding the timezone and keeping the exchange wall-clock date.
They are not converted to UTC before the date is taken.
This matters because \texttt{\^{}VIX} is timestamped in America/Chicago and the other series in America/New\_York, and because daylight-saving offsets change inside a series.

The stage-1 price audit uses \path{data/qqq_ohlc_v1.csv}: \(6{,}967\) rows, 10~March~1999 through 17~November~2026, SHA-256
\begin{center}
\path{963c176144439ec6a6bc31b76c4b045799685d6da9f05ded0c0f3393bd63adbe}.
\end{center}
That file retains raw open, high, low, and close, the vendor adjusted close, and auto-adjusted OHLC.
Stage~2 prices for all six assets are the adjusted panel \path{data/asset_ohlc_main_convention_v1.csv} and the raw panel \path{data/asset_ohlc_raw_v1.csv}.
Both panels have \(14{,}191\) rows on the union of the six calendars, beginning 5~January~1971 because of \texttt{DX-Y.NYB}.
Per-asset calculations drop rows that lack a finite open, high, low, and close rather than feeding the padded union to the indicator library.

\section{Price convention and consistency of the stored indicators}
\label{sec:prices}

The adjusted download calls \texttt{Ticker.history(period='max', interval='1d')} and does not pass \texttt{auto\_adjust}.
In the installed yfinance that argument defaults to true, which is the call made by the research fetcher.
The first downloaded row is dropped together with the missing log return, matching that fetcher.
The raw panel passes \texttt{auto\_adjust=False} and is used only to show that unadjusted equity closes do not reproduce the stored indicators.
The recorded median absolute gaps between adjusted and raw closes are \$31.32 (SPY), \$6.94 (QQQ), \$13.89 (IWM), \$33.68 (TLT), and zero for \texttt{\^{}VIX} and \texttt{DX-Y.NYB}.

The stage-2 fit uses the stored indicator values.
A TA-Lib recomputation, with that library's default arguments, is a consistency check and is not substituted for \(X\).
Write \(x\) for a stored series and \(\hat x\) for the recomputation on adjusted OHLC.
The check requires, for every asset and indicator,
\[
\operatorname{Corr}(x,\hat x)\ge 0.999,
\qquad
\operatorname{median}_t\frac{|x_t-\hat x_t|}{\max(|x_t|,10^{-6})}\le 10^{-3},
\]
a finite positive adjusted close on every feature date, and agreement between the assigned unit family and the empirical rule in Section~\ref{sec:transform}.
The check passed.
The weakest correlation with the adjusted recomputation is \(0.999208\).
The adjusted-panel SHA-256 is
\begin{center}
\path{7e9317cbc535e850cdb34c4aae87cb7f8a826b9ea948b415d8ab29e62c309a6c},
\end{center}
and the raw-panel SHA-256 is
\begin{center}
\path{f486dfa34996f0b1e18a7df5307f4c87cab7f0155cd9cbcd8e9fb5820bd35d8f}.
\end{center}
Both are stored in the stage-2 model metadata, whose creation time is 2026-11-17T21:17:22Z.
The stage-1 metadata records the same feature hash, the QQQ file hash above, and creation time 2026-11-17T20:52:14Z.

\section{Feature preprocessing}
\label{sec:transform}

\subsection{Stage 1}

Stage~1 leaves every stored column unchanged.
Its model directory is \texttt{stage1\_train2018\_v1}.

\subsection{Price-relative transformation}

Stage~2 changes \(132\) columns and leaves \(96\) unchanged.
Let \(C_{a,t}>0\) be the adjusted close of asset \(a\) on the same date \(t\) as the feature.
For a price-level series \(L_{a,t}\) and a price-width series \(V_{a,t}\),
\begin{equation}
\label{eq:ratios}
Y^{\mathrm{level}}_{a,t}=\frac{L_{a,t}}{C_{a,t}}-1,
\qquad
Y^{\mathrm{width}}_{a,t}=\frac{V_{a,t}}{C_{a,t}}.
\end{equation}
Scale-free columns are copied.
Any non-finite result stops the fit.
The formulas and the column list are stored in the stage-2 \path{metadata.json}.

The level family has \(14\) indicators and therefore \(84\) columns: exponential, Kaufman, simple, triple-exponential, and weighted moving averages; linear regression and the linear-regression indicator; median price, typical price, and weighted close; parabolic SAR; and the three Bollinger bands.
The width family has \(8\) indicators and \(48\) columns: average true range, true range, standard deviation, high minus low, MACD, the MACD histogram, ten-day momentum, and the moving-average oscillator.
The unchanged family has \(16\) indicators and \(96\) columns, including the percentage-price oscillator, the high/low ratio, rate of change, TRIX, the bounded oscillators, and the commodity channel index.
Volume indicators are absent because \texttt{\^{}VIX} and \texttt{DX-Y.NYB} do not supply a usable volume series.

\subsection{Why the ratios, and what they do not do}

Stage~1 showed that columns denominated in price leave the 2002--2018 training distribution after 2018.
On that fit the first principal-component score has test mean \(39.02\), and the largest training-standardized column shift is the QQQ average true range, with test mean \(9.24\).
Bollinger bands move with the close, so dividing a band by a range would have treated a price level as a width.
The ratios in \eqref{eq:ratios} remove the nominal price unit: a level becomes a fractional gap from the same-day close, and a width becomes a fraction of that close.

They do not imply stationarity.
A moving average of a trending price remains a smooth, highly persistent series after \eqref{eq:ratios}, and several such averages of the same close remain nearly collinear.
The leading correlation eigenvalue accordingly stays near \(61\), moving only from \(62.03\) to \(61.14\), while the participation ratio moves from \(7.87\) to \(9.60\).
Volatility differences also remain.
The largest absolute column-mean shift after the ratios is the dollar-index average true range, at \(-0.75\) training standard deviations.
The first-component test mean does fall from \(39.02\) to \(-0.38\), which is evidence about the price-era location shift and not evidence of forecast value.

The same-day close is known to a decision made after the close of \(t\).
It would be look-ahead for a decision made at the open of \(t\).
No execution rule is implemented, so this distinction is recorded rather than resolved by a trading convention.

\subsection{Assignment of units}

Families were assigned from the QQQ series before the transformation was accepted.
Let \(P_t\) be the QQQ adjusted close and \(R_t=H_t-L_t\) the same-day dollar range.
A column is a level when its correlation with \(P\) is at least \(0.98\) and its median ratio to \(P\) exceeds \(1/2\).
Otherwise it is a width when its correlation with \(R\) is at least \(0.80\) and exceeds its correlation with \(R/P\), or its correlation with the ten-day dollar change \(\Delta_{10}P\) is at least \(0.80\) and exceeds its correlation with the ten-day percent change, or that dollar-change correlation exceeds the percent-change correlation by more than \(0.15\), is itself above \(0.40\), and the median of \(|x|/P\) is below \(0.05\).
Every other column is scale-free.
Under these rules the Bollinger bands have correlation about \(1\) with the close and a median ratio about \(1.01\), so they are levels.
Standard deviation has correlation \(0.77\) with the close and \(0.82\) with the dollar range, so it is a width.
The empirical family matched the assigned family for all \(38\) indicators.

A different rule was computed and rejected.
The cross-asset slope of median absolute feature value against median price does not separate widths from oscillators: average true range had a slope of about \(-0.12\) because its dollar values are similar across assets whose prices are not.
That slope was not used to assign \eqref{eq:ratios}.

\section{Standardization}
\label{sec:scale}

Both fits use the same training-only standardization.
Let \(X_{\mathrm{tr}}\) be the \(n_{\mathrm{tr}}\times p\) training block of the matrix being decomposed, with \(n_{\mathrm{tr}}=4{,}044\).
The training standard deviation of column \(j\) is the population quantity
\begin{equation}
\label{eq:std}
s_j=\left(\frac{1}{n_{\mathrm{tr}}}\sum_{i=1}^{n_{\mathrm{tr}}}(X_{\mathrm{tr},ij}-\bar X_{\mathrm{tr},j})^2\right)^{1/2}.
\end{equation}
A column with \(s_j\le 10^{-12}\) is dropped.
Neither fit dropped a column.
The retained block is passed to \texttt{StandardScaler}, which subtracts the training mean and divides by \eqref{eq:std}.
The resulting matrix \(Z_{\mathrm{tr}}\) has column means zero and column variances one under the \(1/n_{\mathrm{tr}}\) divisor.
The scaler records \(n_{\mathrm{samples\_seen\_}}=4{,}044\).

Test rows are transformed with those training means and standard deviations.
They are not used to recompute either.
The \(6{,}025\times 228\) matrix produced by applying this frozen map to every row is stored at \path{AI-trades/artifacts/stage1_price_relative_standardized_v1/feature_matrix_standardized.csv}.
The same bytes are copied to \path{AI-trades/artifacts/stage1_price_relative_standardized_v2/feature_matrix_standardized.csv}.
The portable loader is the v2 pipeline described in Section~\ref{sec:roadmap}.
Full-sample standardization was not performed.
Standardization equalizes training variance across assets and indicators.
It does not keep a later price level inside the training distribution; that is the shift measured in Section~\ref{sec:numbers}.

\section{Principal components and the two variance conventions}
\label{sec:pca}

\subsection{The fitted decomposition}

On \(Z_{\mathrm{tr}}\), a full singular-value decomposition is fit with scikit-learn's \texttt{PCA(svd\_solver='full')}.
If \(Z_{\mathrm{tr}}=U\Sigma W^{\top}\) with singular values \(\sigma_1\ge\cdots\ge\sigma_p\), the principal-component scores of a standardized row \(z\) are
\[
f=W^{\top}z.
\]
Test scores use the training \(W\) and the training standardization.
No component is dropped before this diagnostic fit.

\subsection{The \texorpdfstring{\(1/n\)}{1/n} correlation matrix}

Marchenko--Pastur calculations use
\begin{equation}
\label{eq:sn}
S=\frac{1}{n_{\mathrm{tr}}}Z_{\mathrm{tr}}^{\top}Z_{\mathrm{tr}}.
\end{equation}
Because each column of \(Z_{\mathrm{tr}}\) has \(1/n_{\mathrm{tr}}\) variance one, \(\tr(S)=p\) and the mean eigenvalue of \(S\) is \(1\).
The eigenvalues \(\lambda_1\ge\cdots\ge\lambda_p\) of \(S\) are computed with a symmetric eigensolver.
Negative values of magnitude about \(10^{-15}\), which are numerical zeros, are clipped to zero.
These \(\lambda_k\) satisfy
\begin{equation}
\label{eq:singular}
\lambda_k=\frac{\sigma_k^2}{n_{\mathrm{tr}}}.
\end{equation}

\subsection{Scikit-learn's \texorpdfstring{\(1/(n-1)\)}{1/(n-1)} explained variance}

Scikit-learn does not report \(\lambda_k\).
Its \texttt{explained\_variance\_} attribute scales squared singular values by the unbiased divisor,
\begin{equation}
\label{eq:sklearn}
\hat\lambda_k=\frac{\sigma_k^2}{n_{\mathrm{tr}}-1}.
\end{equation}
Combining \eqref{eq:singular} and \eqref{eq:sklearn} gives the exact conversion used in the code,
\begin{equation}
\label{eq:convert}
\lambda_k=\hat\lambda_k\cdot\frac{n_{\mathrm{tr}}-1}{n_{\mathrm{tr}}}.
\end{equation}
The two vectors are compared after this rescaling.
On the stage-1 spectrum their maximum absolute difference is \(9.237\times 10^{-14}\), inside the stored tolerance of \(10^{-6}\).
The \(1/(n-1)\) numbers are not passed to the Marchenko--Pastur edge.
Using them there would inflate every eigenvalue by \(n_{\mathrm{tr}}/(n_{\mathrm{tr}}-1)\) and would make the mean eigenvalue \(n_{\mathrm{tr}}/(n_{\mathrm{tr}}-1)\) rather than \(1\).

The participation ratio, or effective rank, is computed on the \(1/n_{\mathrm{tr}}\) spectrum:
\begin{equation}
\label{eq:pr}
\operatorname{PR}(\lambda)=\frac{\bigl(\sum_{k=1}^{p}\lambda_k\bigr)^2}{\sum_{k=1}^{p}\lambda_k^2}.
\end{equation}
It equals \(1\) when one component carries all of the trace and equals \(p\) when every eigenvalue is equal.

\section{Marchenko--Pastur diagnostics}
\label{sec:mp}

\subsection{Classical edge}

Let \(\gamma=p/n_{\mathrm{tr}}\).
Here \(p=228\) and \(n_{\mathrm{tr}}=4{,}044\), so
\[
\gamma=\frac{228}{4044}=0.056380.
\]
For an \(n\times p\) matrix with i.i.d.\ entries of mean zero and variance \(\sigma^2\), the Marchenko--Pastur theorem places the limiting bulk of \(S\) on
\begin{equation}
\label{eq:mp}
\bigl[\lambda_-,\lambda_+\bigr]
=\Bigl[\sigma^2\bigl(1-\sqrt{\gamma}\bigr)^2,\;
\sigma^2\bigl(1+\sqrt{\gamma}\bigr)^2\Bigr],
\end{equation}
when \(p,n\to\infty\) with \(p/n\to\gamma\in(0,1]\).
At the reference value \(\sigma^2=1\) suggested by column standardization,
\[
\lambda_-=0.581491,\qquad\lambda_+=1.531269.
\]
The diagnostic count is the number of eigenvalues strictly above \(\lambda_+\).
The code does not fit the Marchenko--Pastur density, does not estimate a spiked covariance model, and does not treat an exceedance as a selected factor.

\subsection{Rules that are reported and not applied}

The following counts are written side by side.
Both model files set \texttt{retention\_rule\_selected} to null.

\begin{enumerate}
\item Kaiser: \(\lambda_k>1\).
On the scale of \(S\) this is ``above the mean eigenvalue,'' not a significance test.
\item Marchenko--Pastur at \(\sigma^2=1\): \(\lambda_k>\lambda_+\) from \eqref{eq:mp}.
\item Cumulative variance: the smallest \(k\) such that \(\sum_{i=1}^{k}\lambda_i\big/\sum_{i=1}^{p}\lambda_i\) reaches \(0.80\), \(0.90\), or \(0.95\).
\item An iterated bulk mean, defined next.
It is reported and is not a cutoff, because it does not converge on either fitted spectrum.
\end{enumerate}

The iteration starts at \(\sigma^2_{(0)}=1\).
At step \(m\) it forms the upper edge in \eqref{eq:mp} from \(\sigma^2_{(m-1)}\), averages the eigenvalues at or below that edge, and replaces \(\sigma^2\).
It stops when successive values differ by at most \(10^{-8}\), when the bulk is empty, when the bulk mean is at most \(10^{-8}\), or after \(50\) iterations.
On these spectra the bulk mean collapses toward zero: many correlation eigenvalues sit at numerical zero because the columns are collinear.
The implementation returns a null component count in that case rather than clipping the estimate to a positive constant.

\subsection{Why the classical law does not apply to these indicators}
\label{sec:mp-limits}

Equation \eqref{eq:mp} is a statement about a Wishart-type matrix.
The indicator matrix violates the hypotheses in several ways that the diagnostic does not correct.

The rows are not independent.
Moving averages, oscillators, and ranges at date \(t\) reuse prices from a window ending at \(t\), so successive rows share most of their inputs.
The nominal \(n_{\mathrm{tr}}=4{,}044\) overstates the number of independent covariance observations, and a dependent sample typically spreads the empirical spectrum beyond the i.i.d.\ edge.
An eigenvalue above \(1.531269\) is therefore not identified as a population factor.

The columns are not idiosyncratic noise around a low-rank factor model.
Many are smooth functions of the same adjusted close.
Their large leading eigenvalues measure that constructed collinearity.
Price-relative transformation reduces the common price level and leaves the collinearity of the moving-average family, which is why \(\lambda_1\) barely changes.

Variances are not constant across 2002--2018, and the same matrix is used to estimate the \(p\) means and \(p\) standard deviations that define \(Z_{\mathrm{tr}}\).
The implemented edge is the covariance edge at \(\sigma^2=1\).
It is not the correction appropriate to a sample correlation matrix whose means and variances were estimated in-sample.
No effective-sample-size adjustment, tapering, or random-matrix shrinkage is applied.

The iterated estimator is additionally unstable when the spectrum has a point mass near zero.
Excluding eigenvalues above the current edge lowers the estimated noise variance, which lowers the edge, until the average is dominated by those near-zero eigenvalues.
That path is a numerical diagnostic of the spectrum, not an estimate of \(\sigma^2\).

\section{Training and test separation}
\label{sec:split}

The cutoff \(T=\) 31~December~2018 is inclusive and was fixed before either fit as a first frozen split.
It was not chosen by test-period profitability, and no profitability threshold is defined.
The training sample is the \(4{,}044\) rows from 4~December~2002 through 31~December~2018.
The test sample is the \(1{,}981\) rows from 2~January~2019 through 17~November~2026.
The two blocks are disjoint.
Calendar gaps in the feature index were audited; gaps of five days or less, which cover ordinary weekends, are not treated as missing sessions.

A later run loads the saved scaler and loadings.
If the live column set differs, or the training-row count or training column means differ, the run raises and does not write a new fit over the existing directory.
Stage~2 has its own directory and does not write \path{models/stage1_train2018_v1/}.
Reports and plots are written only when the corresponding diagnostic file is absent, so those artifacts stay with the fit that produced them.
The stage-1 diagnostic Markdown was not regenerated after the reader switched from the snapshot path to the research CSV; its first sample bullet still names the snapshot.
The feature hash of that fit is the hash in Section~\ref{sec:data}.

One cutoff cannot represent a change of dependence inside 2002--2018 or after 2018.
No walk-forward refit is implemented.
A later walk-forward would require a new model directory for each cutoff.

\section{Diagnostic numerical summary}
\label{sec:numbers}

Full tables, scree plots, and column-shift lists are in \path{reports/stage1_diagnostics_v1.md} and \path{reports/stage2_diagnostics_v1.md}, with the eigenvalue and retention CSV files beside them.
The comparison below uses the same split and the same feature hash.
The stage-1 column is the raw matrix.
The stage-2 column is the matrix after \eqref{eq:ratios}.
No row is a selected model.

\begin{center}
\begin{tabular}{lrr}
\toprule
 & Stage 1 & Stage 2 \\
\midrule
Largest eigenvalue \(\lambda_1\) & 62.03 & 61.14 \\
Participation ratio \eqref{eq:pr} & 7.87 & 9.60 \\
PC1 test mean & 39.02 & \(-0.38\) \\
PC1 training standard deviation & 7.88 & 7.82 \\
Largest \(|\)column mean \(z|\) & 9.24 & 0.75 \\
\bottomrule
\end{tabular}
\end{center}

Stage-1 retention counts, none selected, are \(12\), \(19\), \(24\), \(25\), and \(40\) components for cumulative variance \(80\%\), the unit-noise Marchenko--Pastur edge, cumulative variance \(90\%\), Kaiser, and cumulative variance \(95\%\).
The corresponding stage-2 counts are \(17\), \(23\), \(31\), \(30\), and \(49\).
In stage~1 the unit-noise edge retains variance share \(0.8742\); in stage~2 it retains \(0.8585\).
The iterated noise estimate converges in neither fit.

\section{Assumptions and limitations}

A stored feature dated \(t\) is treated as known at the close of \(t\).
The experiment does not rebuild the research indicators except for the stage-2 consistency check, so that check does not prove that every upstream indicator is causal.
The stage-2 divisor is the adjusted close of date \(t\), with the look-ahead distinction in Section~\ref{sec:transform}.

The \(1/n\) and \(1/(n-1)\) conventions answer different questions.
Equation \eqref{eq:sn} is the matrix whose mean eigenvalue is one and whose bulk edge is \eqref{eq:mp}.
Equation \eqref{eq:sklearn} is the variance scikit-learn attaches to a principal-component score.
They differ by the factor in \eqref{eq:convert}; only the former enters the random-matrix diagnostic.

Classical Marchenko--Pastur reasoning does not validate the components above \(\lambda_+\), for the dependence, collinearity, heteroskedasticity, and in-sample standardization reasons in Section~\ref{sec:mp-limits}.
The price-relative map removes a unit.
It does not produce a stationary series, does not orthogonalize the moving-average family, and does not remove volatility shifts.
Overlapping multi-day returns are not computed, so no adjustment for overlapping observations is implemented.
Test-period return is not a fitting or selection criterion.

\section{Methods not implemented}

The following are the unfinished Stage~2 score-matrix export, together with all of Stages~3--7 in Section~\ref{sec:roadmap}.
Roadmap Stage~1 is approved.
Its standardized matrix and portable preprocessor are the v2 artifact in that section.
Nothing in this list is a saved clustering model, a selected dimension, or a completed trading result.

\begin{itemize}
\item A chosen principal-component dimension, including any automatic use of Marchenko--Pastur outliers.
\item Whitening before a distance calculation.
\item \(K\)-means or any other clustering, and any search over \(K\).
\item A Markov transition matrix, and five-day or ten-day regime probabilities.
\item Conditional QQQ returns, a buy-and-hold comparison, transaction costs, and a trading threshold.
\item Paper-trading logs, orders, or a daily decision rule.
\item Walk-forward refits.
\end{itemize}

Clustering and trading require an approval separate from this methods note.
Writing or compiling this document does not supply that approval.

\end{document}
